\ifdefined\XeTeXversion\else\ifdefined\directlua\else\pdfoutput=1\fi\fi % arXiv: pdfLaTeX (keep in the first 5 lines)
\documentclass[11pt]{article}
\usepackage{thefinai}
% Options: [numbers] numeric citations · [anonymous] blind review · [review] line numbers · [nologo]

% ------------------------------------------------------------------ title card
\title{FinBench-X: A Template Title for a Financial AI Benchmark Paper}
\date{October 10, 2026}
% \reportlabel{Technical Report}   % optional small label above the title

\author{%
  First Author\textsuperscript{1,*}, Second Author\textsuperscript{1,2,*}, Third Author\textsuperscript{3},
  Fourth Author\textsuperscript{2}, Fifth Author\textsuperscript{1,\dag}%
}
\affiliation[1]{The Fin AI}
\affiliation[2]{University One}
\affiliation[3]{Institute Two}
\authornote{\textsuperscript{*}Equal contribution}
\authornote{\textsuperscript{\dag}Corresponding author}

\correspondence{\href{mailto:contact@thefin.ai}{contact@thefin.ai}}
\metadata[Code]{\ghlink{The-FinAI/FinBen}}
\metadata[Hugging Face]{\hflink{collections/TheFinAI}}
\metadata[Leaderboard]{\href{https://huggingface.co/spaces/TheFinAI/Open-FinLLM-Leaderboard}{Open FinLLM Leaderboard}}

\begin{document}

% The abstract goes before \maketitle so it is placed inside the title card.
\begin{abstract}
Replace this paragraph with a 150--250 word abstract. State the problem in financial AI you address, what you
release (a benchmark, model, dataset or agent system), how you evaluate it, and the one or two results a reader
should remember. Keep numbers concrete and claims checkable. The date, correspondence and links below the abstract
come from \texttt{\textbackslash date}, \texttt{\textbackslash correspondence} and \texttt{\textbackslash metadata}.
\end{abstract}

\maketitle

% A teaser figure on page 1, as most current reports do (illustrative numbers).
\begin{figure}[h]
  \centering
  \begin{tikzpicture}
    \begin{axis}[finai,finaibars,width=\linewidth,height=4.6cm,bar width=9pt,ymin=0,ymax=90,
      symbolic x coords={Sentiment,QA,Extraction,Forecasting,Trading},xtick=data,
      ylabel={Score},enlarge x limits=0.12,legend columns=-1,legend style={at={(0.5,1.04)},anchor=south},
      nodes near coords,every node near coord/.append style={font=\sffamily\tiny,color=finaisoft}]
      \addplot coordinates {(Sentiment,82) (QA,68) (Extraction,58) (Forecasting,54) (Trading,61)};
      \addplot coordinates {(Sentiment,78) (QA,63) (Extraction,55) (Forecasting,49) (Trading,52)};
      \addplot coordinates {(Sentiment,71) (QA,57) (Extraction,47) (Forecasting,45) (Trading,44)};
      \legend{Ours (8B), Model A (70B), Model B (8B)}
    \end{axis}
  \end{tikzpicture}
  \caption{A teaser figure summarizing the main result. Numbers are illustrative.}
  \label{fig:teaser}
\end{figure}

% ------------------------------------------------------------------ body
\section{Introduction}
Large language models are increasingly used for financial analysis, reporting and decision making
\citep{xie2023pixiu,xie2024finben}. Introduce the gap, then list the contributions:

\begin{itemize}
  \item \textbf{Benchmark.} What is new about the tasks, data sources or evaluation protocol.
  \item \textbf{Evaluation.} Which models were tested, and under which settings.
  \item \textbf{Findings.} The main empirical takeaways (\cref{sec:results}).
\end{itemize}

\begin{takeaway}
Use this box for the one or two sentences a busy reader should remember.
\end{takeaway}

\section{Related Work}
Position the work against financial benchmarks \citep{xie2024finben,peng2025multifinben} and reasoning
models \citep{qian2025fino1}. Group prior work by theme rather than listing papers one by one.

\section{Benchmark Construction}
\label{sec:method}
\subsection{Task Design}
Describe the tasks, label space and data sources. Define notation where needed, for example accuracy over a test
set $\mathcal{D}$:
\begin{equation}
  \mathrm{Acc} = \frac{1}{|\mathcal{D}|}\sum_{(x,y)\in\mathcal{D}} \mathbb{I}\left[f(x) = y\right].
\end{equation}

\subsection{Prompting}
\begin{promptbox}[Zero-shot prompt (sentiment)]
Analyze the sentiment of the following financial news headline. Answer with one word:
positive, negative or neutral.\\
Headline: \{text\}
\end{promptbox}

\section{Results}
\label{sec:results}
\Cref{tab:main} shows the table style: booktabs rules, \best{bold} for the best and \second{underline} for the
second-best result, and a tinted row for your method (\texttt{\textbackslash oursrow}).

\begin{table}[t]
  \caption{Main results (illustrative numbers). Higher is better.}
  \label{tab:main}
  \centering\small
  \begin{tabular}{lccccc}
    \toprule
    \tablehead{Model} & \tablehead{Sentiment} & \tablehead{QA} & \tablehead{Extraction} & \tablehead{Trading} & \tablehead{Avg.} \\
    \midrule
    Model A (70B) & \second{78.0} & \second{63.0} & \second{55.0} & \second{52.0} & \second{62.0} \\
    Model B (8B)  & 71.0 & 57.0 & 47.0 & 44.0 & 54.8 \\
    \oursrow Ours (8B) & \best{82.0} & \best{68.0} & \best{58.0} & \best{61.0} & \best{67.3}\up{5.3} \\
    \bottomrule
  \end{tabular}
\end{table}

\begin{finding}{Finding 1.}
State each empirical result in one or two sentences, then point to the table or figure that supports it.
\end{finding}

\section{Conclusion}
Summarize what was released, what it shows, and the main limitations.

\section*{Limitations and Ethics}
Financial data and model outputs can be wrong; nothing in this work is investment advice. Describe data licenses,
personal-data handling and intended use.

\section*{Acknowledgments}
Thank funders, compute providers and annotators.

\bibliographystyle{plainnat}
\bibliography{references}

% ------------------------------------------------------------------ appendix
\appendix
\section{Author Contributions}
List each author's role (conceptualization, data, experiments, writing).

\section{Additional Results}
Extra tables, prompts and examples go here.

\end{document}
